Fix \(q\ge 1\), put \(n=4q\), and fix \(a,b,c_j,s_j>0\). These positivity conditions are the only parameter boundary conditions used below. The supports of distinct block columns of \(U_j\) are disjoint, and every column has positive squared norm, so \(U_j\) has full column rank \(q\). Thus the response ellipsoids and all eigenvalues in the statement are well defined. Boundary designs assigning zero probability to some local directions are allowed by \(\mathrm{def:design-space}\); the proof uses no inverse exposure probability or overlap condition.

For compactness, define the three squared block loadings
\[
d_{1,1}=c_1,\qquad d_{1,2}=c_1+s_1a,\qquad d_{1,3}=c_1+s_1b,
\]
\[
d_{2,1}=c_2+s_2a,\qquad d_{2,2}=c_2,\qquad d_{2,3}=c_2+s_2b. \tag{1}
\]
The identities \(h_k=v_k/2\) and \(h_k^\top h_{k'}=\mathbf 1\{k=k'\}\) imply that, on block \(\ell\),
\[
Z_{[\ell]}=\varepsilon v_k
\quad\Longrightarrow\quad
u_{j\ell}^\top Z=2\varepsilon\sqrt{d_{j,k}},
\qquad
(u_{j\ell}^\top Z)^2=4d_{j,k}. \tag{2}
\]
This calculation is pointwise in the assignment and does not require independence across blocks. Since \(C_j(\pi)=\mathbb E_\pi[(U_j^\top Z)(U_j^\top Z)^\top]\), summing its diagonal entries and applying \(\mathrm{def:direction-frequencies}\) yields
\[
\frac{\operatorname{tr}C_1(\pi)}q
=4\{c_1+s_1(ap_2(\pi)+bp_3(\pi))\}, \tag{3}
\]
\[
\frac{\operatorname{tr}C_2(\pi)}q
=4\{c_2+s_2(ap_1(\pi)+bp_3(\pi))\}. \tag{4}
\]
Each \(C_j(\pi)\) is a positive-semidefinite \(q\)-by-\(q\) matrix. Its largest eigenvalue is at least the average of its \(q\) nonnegative eigenvalues, so
\[
\lambda_{\max}\{C_j(\pi)\}\ge q^{-1}\operatorname{tr}C_j(\pi). \tag{5}
\]
Equations (3)--(5) prove the two spectral certificates.

Because all frequencies and increments are nonnegative, (3)--(5) and \(\mathrm{lem:variance-and-spectral-risk}\) imply \(R_j(\pi)\ge 4c_j/n^2\). For class \(1\), choose direction \(v_1\) in every block and multiply it by independent symmetric signs. Then (2) gives independent centered scores of squared magnitude \(4c_1\), and hence \(C_1=4c_1I_q\). The analogous product law on \(v_2\) gives \(C_2=4c_2I_q\). Both laws are blinded members of \(\Pi_n\). Therefore
\[
r_j^\star=\frac{4c_j}{n^2}. \tag{6}
\]
Under block complete randomization, every direction has probability \(1/3\), and independent symmetric block signs make the score covariance scalar. Equations (1)--(2) give
\[
C_j(\pi_0)=4\left\{c_j+\frac{s_j(a+b)}3\right\}I_q. \tag{7}
\]
Combining (6)--(7) with the spectral-risk identity proves
\[
B_j=\frac4{n^2}\left\{c_j+\frac{s_j(a+b)}3\right\},
\qquad
\Delta_j=\frac{4s_j(a+b)}{3n^2}>0. \tag{8}
\]
The strict inequality in (8) follows from \(s_j,a,b>0\), so every division by \(\Delta_j\) below is legitimate.

It remains to solve the design problem. Suppress the argument \(\pi\), set \(A=a+b\), and define
\[
L_1=ap_2+bp_3,\qquad L_2=ap_1+bp_3. \tag{9}
\]
The trace certificates, (8), and the definition of \(g_j\) imply
\[
g_1(\pi)\le 1-\frac{3L_1}{A},
\qquad
g_2(\pi)\le 1-\frac{3L_2}{A}. \tag{10}
\]
Each block uses exactly one of the three direction pairs, so \(p_1+p_2+p_3=1\). Hence
\[
\max(L_1,L_2)
\ge \frac{L_1+L_2}{2}
=\frac a2(1-p_3)+bp_3
\ge \min\left\{\frac a2,b\right\}. \tag{11}
\]
Consequently every feasible design satisfies
\[
\min_j g_j(\pi)
\le 1-\frac{3}{a+b}\min\left\{\frac a2,b\right\}
=
\begin{cases}
\dfrac{a-2b}{a+b},&a\ge 2b,\\[6pt]
\dfrac{2b-a}{2(a+b)},&a\le 2b.
\end{cases} \tag{12}
\]

The certificate is attainable. Consider a product law which independently in each block chooses direction \(k\) with probability \(p_k\) and then applies an independent symmetric sign. The score coordinates are centered and independent across blocks, while (2) gives identical second moments. Therefore
\[
C_1=4(c_1+s_1L_1)I_q,
\qquad
C_2=4(c_2+s_2L_2)I_q. \tag{13}
\]
If \(a>2b\), take \((p_1,p_2,p_3)=(0,0,1)\); then \(L_1=L_2=b\). If \(a<2b\), take \((p_1,p_2,p_3)=(1/2,1/2,0)\); then \(L_1=L_2=a/2\). If \(a=2b\), take any \(p_1=p_2=t\) and \(p_3=1-2t\) with \(0\le t\le 1/2\); then again \(L_1=L_2=b=a/2\). In all cases the bounds (10)--(12) are equalities by (13). Thus the upper bound (12) is attained and proves the displayed piecewise formula for \(\lambda^\star\). Equation (13), together with the trace lower bound, is the promised matching finite-sample optimality certificate. The first witness is exactly \(\pi_{\mathrm{sh}}\).

We next characterize the optimal average frequencies. If a design is minimax, equality must hold throughout the chain leading from (10) through (11) to (12). When \(a>2b\), the affine average in (11) equals
\[
b+\left(\frac a2-b\right)(1-p_3),
\]
so equality forces \(p_3=1\), and therefore \(p_1=p_2=0\). When \(a<2b\), equality forces \(p_3=0\); equality between the maximum and the average then forces \(L_1=L_2\), hence \(p_1=p_2=1/2\). When \(a=2b\), the affine average in (11) is identically \(b\), and equality of the maximum and the average is equivalent to \(L_1=L_2\), or \(p_1=p_2=t\) with \(p_3=1-2t\) and \(0\le t\le 1/2\). Conversely, the product laws used in (13) attain the frontier for every displayed frequency vector. This proves the exact frequency characterization.

For a regime-specific product witness, let \(m=\min\{a/2,b\}\). Equations (7) and (13) give
\[
B_j-R_j
=\frac{4s_j}{n^2}\left\{\frac{a+b}{3}-m\right\}
=
\begin{cases}
\dfrac{4s_j(a-2b)}{3n^2},&a>2b,\\[6pt]
\dfrac{2s_j(2b-a)}{3n^2},&a<2b,\\[6pt]
0,&a=2b.
\end{cases} \tag{14}
\]
Thus both class risks improve strictly exactly off the elbow. At \(a=2b\), block complete randomization has \(p_1=p_2=p_3=1/3\), belongs to the optimal frequency family, and its scalar matrices in (7) attain (12); hence it is maximin and \(\lambda^\star=0\).

For the required smallest-instance check, take \(q=1\) and \(n=4\). Every \(C_j\) is scalar, and a blinded design is described by its masses \((p_1,p_2,p_3)\) on the three sign pairs. Direct enumeration of the six assignments gives exactly (3)--(4), so minimizing the larger of the two scalar increments is precisely (11). The three equality cases above exhaust the simplex. This verifies the frontier, its elbow, and all boundary-frequency claims without an asymptotic approximation.